\documentclass[11pt]{article}
\usepackage[a4paper,margin=18mm]{geometry}
\usepackage[no-math]{fontspec}
\setmainfont{Latin Modern Roman}
\usepackage{amsmath,amssymb,amsthm,xfrac,xspace,hyperref,graphicx,comment}
\excludecomment{DefTralics}

\providecommand{\bibliofont}{\small}
\newcommand{\ie}{i.e.,\xspace}
\newcommand{\eg}{e.g.,\xspace}
\newcommand{\etc}{etc.\xspace}
\newcommand{\cf}{cf.\xspace}
\newcommand{\resp}{resp.\xspace}
\newcommand{\smaller}{\small}
\newcommand{\Subsection}{\subsection}
\newcommand{\Subsubsection}{\subsubsection}
\newcommand{\skpt}{\smallskip}
\emergencystretch=3em
\numberwithin{equation}{section}\renewcommand{\encodingdefault}{TU}\input{author-preamble.tex}
\renewcommand{\encodingdefault}{TU}\AtBeginDocument{\fontencoding{TU}\selectfont}
\begin{document}
\begin{center}{\Large Stable item 2050}\end{center}
\paragraph{Source and attribution.}
Frédéric Déglise, Jean Fasel, Fangzhou Jin, Adeel A. Khan, \emph{On the rational motivic homotopy category}, Journal de l’École polytechnique — Mathématiques 8 (2021), publisher version, DOI 10.5802/jep.153. \href{https://jep.centre-mersenne.org/articles/10.5802/jep.153/}{Publisher article}. Authors' text: \href{https://creativecommons.org/licenses/by/4.0/}{CC BY 4.0}.
\paragraph{Editorial problem boundary.}
Prove only Theorem 3.6(2) below for the minus and rational motivic sphere spectra, using the exact smoothable-morphism fundamental-class definition of absolute purity between regular noetherian schemes and the printed quasi-compact/quasi-separated conventions. Full new characteristic-zero reduction, exceptional-pullback extension/joint conservativity and the complete first equal-characteristic Popescu argument are retained. Category-of-Witt-motives identification and separate formal six-operation corollaries are context only.
\paragraph{Difficulty (editorial): H3.}
The cited Popescu and prior relative-purity results do not directly establish the smoothable absolute-purity map over regular bases of mixed characteristic. The author proves the characteristic-zero reduction and develops the required exceptional-pullback extension and joint conservativity through continuity and essentially finite-presentation approximations. The complete first equal-characteristic purity argument then closes the reduction, rather than assuming absolute purity as a general six-operation identity.
\paragraph{Editorial transcription note.}
Original author language, mathematics and ordering are preserved. Publisher front matter is replaced by this independent article wrapper. Bibliography is recovered from matching cached publisher HTML, including its TeX formulas. No new proof or mathematical repair is supplied. Surrounding original results are retained for conservative context and are not extra selected corpus claims.
\clearpage
\input{author-body.tex}
\begin{thebibliography}{99}
\bibitem{SGA4} M. Artin, A. Grothendieck \& J.-L. Verdier - Théorie des topos et cohomologie étale des schémas , Lect. Notes in Math. , vol. 269, 270, 305 , Springer-Verlag, 1972–1973, Séminaire de Géométrie Algébrique du Bois–Marie 1963–64 (SGA 4)
\bibitem{ALP} A. Ananyevskiy, M. Levine \& I. Panin - “Witt sheaves and the $\eta $-inverted sphere spectrum” , J. Topology 10 (2017) no. 2, p. 370-385
\bibitem{AnaSL} A. Ananyevskiy - “$\mathrm{SL}$-oriented cohomology theories” , 2019
\bibitem{AyoubThesis} J. Ayoub - Les six opérations de Grothendieck et le formalisme des cycles évanescents dans le monde motivique , Astérisque , vol. 314-315 , Société Mathématique de France, Paris, 2007
\bibitem{AyoubEt} J. Ayoub - “La réalisation étale et les opérations de Grothendieck” , Ann. Sci. École Norm. Sup. (4) 47 (2014) no. 1, p. 1-145
\bibitem{Bachmann} T. Bachmann - “Motivic and real étale stable homotopy theory” , Compositio Math. 154 (2018) no. 5, p. 883-917
\bibitem{BalmerWitt} P. Balmer - “Witt cohomology, Mayer-Vietoris, homotopy invariance and the Gersten conjecture” , $K$-Theory 23 (2001) no. 1, p. 15-30
\bibitem{BalmerBasic} P. Balmer - “Witt groups” , in Handbook of $K$-theory. Vol. 1, 2 , Springer, Berlin, 2005, p. 539-576
\bibitem{BBD} A. A. Beĭlinson, J. Bernstein \& P. Deligne - “Faisceaux pervers” , in Analysis and topology on singular spaces, I (Luminy, 1981) , Astérisque , vol. 100 , Société Mathématique de France, Paris, 1982, p. 5-171
\bibitem{BCDFO} T. Bachmann, B. Calmès, F. Déglise, J. Fasel \& P. A. Østvær - “Milnor-Witt motives” , 2020
\bibitem{BD1} M. Bondarko \& F. Déglise - “Dimensional homotopy t-structures in motivic homotopy theory” , Adv. Math. 311 (2017), p. 91-189
\bibitem{BachmannFasel} T. Bachmann \& J. Fasel - “On the effectivity of spectra representing motivic cohomology theories” , 2018
\bibitem{BalmerGillePaninWalter} P. Balmer, S. Gille, I. Panin \& C. Walter - “The Gersten conjecture for Witt groups in the equicharacteristic case” , Doc. Math. 7 (2002), p. 203-217
\bibitem{BachmannHoyois} T. Bachmann \& M. Hoyois - Norms in motivic homotopy theory , Astérisque , Société Mathématique de France, Paris, 2021, to appear
\bibitem{BlochOgus} S. Bloch \& A. Ogus - “Gersten’s conjecture and the homology of schemes” , Ann. Sci. École Norm. Sup. (4) 7 (1974) no. 4, p. 181-201
\bibitem{Bondarkorweight} M. Bondarko - “Weights for relative motives: relation with mixed complexes of sheaves” , Internat. Math. Res. Notices (2014) no. 17, p. 4715-4767
\bibitem{BalmerWalter} P. Balmer \& C. Walter - “A Gersten-Witt spectral sequence for regular schemes” , Ann. Sci. École Norm. Sup. (4) 35 (2002) no. 1, p. 127-152
\bibitem{CD5} D.-C. Cisinski \& F. Déglise - “Integral mixed motives in equal characteristics” , Doc. Math. (2015), p. 145-194, Extra volume: Alexander S. Merkurjev’s sixtieth birthday
\bibitem{CDEtale} D.-C. Cisinski \& F. Déglise - “Étale motives” , Compositio Math. 152 (2016) no. 3, p. 556-666
\bibitem{CD3} D.-C. Cisinski \& F. Déglise - Triangulated categories of mixed motives , Springer Monographs in Math. , Springer, Cham, 2019
\bibitem{realK1} B. Calmès, E. Dotto, J. Harpaz, F. Hebestreit, M. Land, K. Moi, D. Nardin, T. Nikolaus \& W. Steimle - “Hermitian $K$-theory for stable $\infty $-categories I: Foundations” , 2020
\bibitem{realK2} B. Calmès, E. Dotto, J. Harpaz, F. Hebestreit, M. Land, K. Moi, D. Nardin, T. Nikolaus \& W. Steimle - “Hermitian $K$-theory for stable $\infty $-categories II: Cobordism categories and additivity” , 2020
\bibitem{realK3} B. Calmès, E. Dotto, J. Harpaz, F. Hebestreit, M. Land, K. Moi, D. Nardin, T. Nikolaus \& W. Steimle - “Hermitian $K$-theory for stable $\infty $-categories III: Grothendieck-Witt groups of rings” , 2020
\bibitem{CF1} B. Calmès \& J. Fasel - “Finite Chow-Witt correspondences” , 2014
\bibitem{CisLect} D.-C. Cisinski - “Cohomological methods in intersection theory” (2019), arXiv:1905.03478
\bibitem{CTHK} J.-L. Colliot-Thélène, R. Hoobler \& B. Kahn - “The Bloch-Ogus-Gabber theorem” , in Algebraic $K$-theory (Toronto, ON, 1996) , Fields Inst. Commun. , vol. 16 , American Mathematical Society, Proovidence, RI, 1997, p. 31-94
\bibitem{SGA4.5} P. Deligne - Cohomologie étale , Lect. Notes in Math. , vol. 569 , Springer-Verlag, 1977, Séminaire de Géométrie Algébrique du Bois–Marie SGA $4\frac{1}{2}$
\bibitem{Del} P. Deligne - “Le déterminant de la cohomologie” , in Current trends in arithmetical algebraic geometry (Arcata, Calif., 1985) , Contemp. Math. , vol. 67 , American Mathematical Society, Providence, RI, 1987, p. 93-177
\bibitem{DF3} F. Déglise \& J. Fasel - “The Borel character” , 2020
\bibitem{DFJKBorel} F. Déglise, J. Fasel, F. Jin \& A. A. Khan - “Borel isomorphism and absolute purity” , 2019
\bibitem{DJK} F. Déglise, F. Jin \& A. A. Khan - “Fundamental classes in motivic homotopy theory” , J. Eur. Math. Soc. (JEMS) (2021), to appear
\bibitem{DegliseBivariant} F. Déglise - “Bivariant theories in motivic stable homotopy” , Doc. Math. 23 (2018), p. 997-1076
\bibitem{DegliseOrientation} F. Déglise - “Orientation theory in arithmetic geometry” , in $K$-Theory—Proceedings of the International Colloquium (Mumbai, 2016) , Hindustan Book Agency, New Delhi, 2018, p. 239-347
\bibitem{EHKSY3} E. Elmanto, M. Hoyois, A. A. Khan, V. Sosnilo \& M. Yakerson - “Modules over algebraic cobordism” , Forum Math. Pi 8 (2020), article ID e14, 44 pages
\bibitem{ElmantoKhan} E. Elmanto \& A. A. Khan - “Perfection in motivic homotopy theory” , Proc. London Math. Soc. (3) 120 (2020) no. 1, p. 28-38
\bibitem{ElmantoKolderup} E. Elmanto \& H. Kolderup - “On modules over motivic ring spectra” , Ann. $K$-Theory 5 (2020) no. 2, p. 327-355
\bibitem{EKM} R. Elman, N. Karpenko \& A. Merkurjev - The algebraic and geometric theory of quadratic forms , AMS Colloquium Publications , vol. 56 , American Mathematical Society, Providence, RI, 2008
\bibitem{Fasel_wCH} J. Fasel - Groupes de Chow-Witt , Mém. Soc. Math. France (N.S.) , vol. 113 , Société Mathématique de France, Paris, 2008
\bibitem{Feld2} N. Feld - “Morel homotopy modules and Milnor-Witt cycle modules” , 2019
\bibitem{Feld1} N. Feld - “Milnor-Witt cycle modules” , J. Pure Appl. Algebra 224 (2020) no. 7, p. 41
\bibitem{FaselSrinivas} J. Fasel \& V. Srinivas - “Chow-Witt groups and Grothendieck-Witt groups of regular schemes” , Adv. Math. 221 (2009) no. 1, p. 302-329
\bibitem{Fuji} K. Fujiwara - “A proof of the absolute purity conjecture (after Gabber)” , in Algebraic geometry 2000, Azumino (Hotaka) , Adv. Stud. Pure Math. , vol. 36 , Math. Soc. Japan, Tokyo, 2002, p. 153-183
\bibitem{Ful} W. Fulton - Intersection theory , Ergeb. Math. Grenzgeb. (3) , vol. 2 , Springer-Verlag, Berlin, 1998
\bibitem{Garkusha} G. Garkusha - “Reconstructing rational stable motivic homotopy theory” , Compositio Math. 155 (2019) no. 7, p. 1424-1443
\bibitem{Gille} S. Gille - “A graded Gersten-Witt complex for schemes with a dualizing complex and the Chow group” , J. Pure Appl. Algebra 208 (2007) no. 2, p. 391-419
\bibitem{EGA4} A. Grothendieck - “Éléments de géométrie algébrique. IV. Étude locale des schémas et des morphismes de schémas. I” , Publ. Math. Inst. Hautes Études Sci. 20 (1964), p. 5-259
\bibitem{SGA5} A. Grothendieck - Cohomologie $\ell $-adique et fonctions ${L}$ , Lect. Notes in Math. , vol. 589 , Springer-Verlag, 1977, Séminaire de Géométrie Algébrique du Bois–Marie 1965–66 (SGA 5)
\bibitem{HartshorneRD} R. Hartshorne - Residues and duality , Lect. Notes in Math. , vol. 20 , Springer-Verlag, Berlin-New York, 1966
\bibitem{HoyoisLefschetz} M. Hoyois - “A quadratic refinement of the Grothendieck-Lefschetz-Verdier trace formula” , Algebraic Geom. Topol. 14 (2014) no. 6, p. 3603-3658
\bibitem{Hebert} D. Hébert - “Structure de poids à la Bondarko sur les motifs de Beilinson” , Compositio Math. 147 (2011) no. 5, p. 1447-1462
\bibitem{Gabber} - Travaux de Gabber sur l’uniformisation locale et la cohomologie étale des schémas quasi-excellents (L. Illusie, Y. Laszlo \& F. Orgogozo, eds.) , Astérisque , vol. 363-364 , Société Mathématique de France, Paris, 2014
\bibitem{Jacobson} J. Jacobson - “Real cohomology and the powers of the fundamental ideal in the Witt ring” , Ann. $K$-Theory 2 (2017) no. 3, p. 357-385
\bibitem{Jin} F. Jin - “Borel–Moore motivic homology and weight structure on mixed motives” , Math. Z. 283 (2016) no. 3, p. 1149-1183
\bibitem{KhanThesis} A. A. Khan - Motivic homotopy theory in derived algebraic geometry , Ph. D. Thesis, Universität Duisburg-Essen, 2016
\bibitem{KhanVirtual} A. A. Khan - “Virtual fundamental classes of derived stacks I” , 2019
\bibitem{KhanSix} A. A. Khan - “Voevodsky’s criterion for constructible categories of coefficients” (2021), Preprint, available at https://www.preschema.com/papers/six.pdf
\bibitem{Kne} M. Knebusch - “Symmetric bilinear forms over algebraic varieties” , in Conference on Quadratic Forms—1976 (Kingston, Ont., 1976) , Queen’s Papers in Pure and Appl. Math. , vol. 46 , 1977, p. 103-283
\bibitem{LamQuadratic} T. Y. Lam - Introduction to quadratic forms over fields , Graduate Studies in Math. , vol. 67 , American Mathematical Society, Providence, RI, 2005
\bibitem{LurieHTT} J. Lurie - Higher topos theory , Annals of Math. Studies , vol. 170 , Princeton University Press, Princeton, NJ, 2009
\bibitem{LurieHA} J. Lurie - “Higher algebra” (2012), Preprint, available at https://www.math.ias.edu/\textasciitilde{}lurie/papers/HigherAlgebra.pdf
\bibitem{LurieSAG} J. Lurie - “Spectral algebraic geometry” (2018), Preprint, available at https://www.math.ias.edu/\textasciitilde{}lurie/papers/SAG-rootfile.pdf
\bibitem{Morelpi0} F. Morel - “On the motivic $\pi _0$ of the sphere spectrum” , in Axiomatic, enriched and motivic homotopy theory , NATO Sci. Ser. II Math. Phys. Chem. , vol. 131 , Kluwer Acad. Publ., 2004, p. 219-260
\bibitem{MorelSplitting} F. Morel - “Rational stable splitting of Grassmannians and rational motivic sphere spectrum” , 2006
\bibitem{MorelA1} F. Morel - $\mathbb{A}^1$-algebraic topology over a field , Lect. Notes in Math. , vol. 2052 , Springer, Heidelberg, 2012
\bibitem{MV} F. Morel \& V. Voevodsky - “${\mathbb{A}}^1$-homotopy theory of schemes” , Publ. Math. Inst. Hautes Études Sci. (1999) no. 90, p. 45-143
\bibitem{Panin} I. Panin - “Homotopy invariance of the sheaf $W_{\mathrm{Nis}}$ and of its cohomology” , in Quadratic forms, linear algebraic groups, and cohomology , Dev. Math. , vol. 18 , Springer, New York, 2010, p. 325-335
\bibitem{PaninWalter1} I. Panin \& C. Walter - “On the motivic commutative ring spectrum BO” , St. Petersburg Math. J. 30 (2019) no. 6, p. 933–972
\bibitem{Rob} M. Robalo - “$K$-theory and the bridge from motives to noncommutative motives” , Adv. Math. 269 (2015), p. 399-550
\bibitem{RO} O. Röndigs \& P. A. Østvær - “On modules over motivic ring spectra” , Adv. Math. 219 (2008) no. 2, p. 689–727
\bibitem{Scheiderer} C. Scheiderer - Real and étale cohomology , Lect. Notes in Math. , vol. 1588 , Springer-Verlag, Berlin, 1994
\bibitem{Scholl} A. Scholl - “Integral elements in $K$-theory and products of modular curves” , in The arithmetic and geometry of algebraic cycles (Banff, AB, 1998) , NATO Sci. Ser. C Math. Phys. Sci. , vol. 548 , Kluwer Acad. Publ., 2000, p. 467-489
\bibitem{Schlicht2} M. Schlichting - “Hermitian $K$-theory, derived equivalences and Karoubi’s fundamental theorem” , J. Pure Appl. Algebra 221 (2017) no. 7, p. 1729-1844
\bibitem{Spivakovsky} M. Spivakovsky - “A new proof of D. Popescu’s theorem on smoothing of ring homomorphisms” , J. Amer. Math. Soc. 12 (1999) no. 2, p. 381-444
\bibitem{SpitzweckSpectrum} M. Spitzweck - A commutative $\mathbb{P}^1$-spectrum representing motivic cohomology over Dedekind domains , Mém. Soc. Math. France (N.S.) , vol. 157 , Société Mathématique de France, Paris, 2018
\bibitem{SchlichTri} M. Schlichting \& G. S. Tripathi - “Geometric models for higher Grothendieck-Witt groups in $\mathbb{A}^1$-homotopy theory” , Math. Ann. 362 (2015) no. 3-4, p. 1143-1167
\bibitem{Stacks} Stacks project authors - “The Stacks project” , https://stacks.math.columbia.edu , 2021
\bibitem{Thomason1} R. W. Thomason - “Absolute cohomological purity” , Bull. Soc. math. France 112 (1984) no. 3, p. 397-406
\bibitem{TT} R. W. Thomason \& T. Trobaugh - “Higher algebraic $K$-theory of schemes and of derived categories” , in The Grothendieck Festschrift, Vol. III , Progress in Math. , vol. 88 , Birkhäuser Boston, Boston, MA, 1990, p. 247-435
\end{thebibliography}
\end{document}
